\subsection{The retained arithmetic interface}

Write \(t=16+d^2\), where \(d\) is an unrestricted integer. The principal
arithmetic quantities are
\[
 F,\quad B,\quad V,\quad x,\quad y,\quad N,\quad q,\quad
 M,\qquad H=1+\sum_{\nu=1}^{9}a_\nu.
\]
The \(a_\nu\) are the nine actual CRT representatives. In the stored source,
several are a coordinate plus one; it is their actual values, rather than
their shifted coordinates, that will become nonnegative.

The retained exact-power requests have the intended meanings
\begin{equation}\label{eq:power-meanings}
 V=(F^2+1)^t,\qquad N=B^x,\qquad q=(BH)^x,\qquad M=q^x,
\end{equation}
together with the powers used in the binomial numerators and denominators.
There are eighteen requests, sharing fourteen independent principal Pell
pairs. At this point these displays are requested meanings, not assumptions.
Their interpretation is recovered in a specified order below.

The beta spacing is left literally unchanged:
\begin{equation}\label{eq:beta}
 y=V(V-1).
\end{equation}
The endpoint equations include
\begin{equation}\label{eq:endpoints}
 x=F+1+(y+1)u_0,\qquad
 x=1+\bigl(1+(t+1)y\bigr)u_t.
\end{equation}
Initially \(u_0,u_t,x,F,B,q\) may have either sign.

The eleven normalized remainder registers fall into four groups:
\begin{center}
\begin{tabular}{llr}
\hline
Denominator & Meaning & Number\\
\hline
\(q^t\) & Factorial digit and encoded index &2\\
\(q^{t+1}\) & Four syntax representatives and derived label &5\\
\(q^{y+1}\) & Two pending-multiset representatives &2\\
\(q^x\) & Two beta quotients &2\\
\hline
\end{tabular}
\end{center}
The existing ratio equations are
\begin{equation}\label{eq:ratios}
 q^t v_t=M,\qquad q^{t+1}v_{t+1}=M,\qquad
 q^{y+1}v_y=M.
\end{equation}
Each displayed power in these equations is a retained output register.
The quotient for the last group is the constant one.

For clarity, put
\[
 G(a,b,c)=a^2+b^2+c(c+1).
\]
This form is nonnegative at every integer triple and represents every
nonnegative integer. The latter is the familiar consequence of the
three-square theorem applied to \(4n+1\): exactly one square is odd, and
the two even squares can be divided by four. This is the same guard form
already used in the parent.

The inherited signed affine-radix core replaces three separate initial
scalar sign certificates by the first three rows below. This round
changes only its last row, from $M^2=1+4(S^2+\sum R_i^2)+G$ to
\begin{align}
 q&=2z+1,                                                    \label{eq:odd}\\
 Bv&=M,                                                     \label{eq:factor-ratio}\\
 MF&=vV+S,                                                  \label{eq:factor-error}\\
 M^2&=2\left(S^2+\sum_{i=1}^{11}R_i^2\right)+G(a,b,c).
                                                               \label{eq:common-norm}
\end{align}
The integers \(v,S\) normalize the factorial error at the same
denominator \(M\) as the other remainders. No inequality, division,
factorial, or variable power has been added to the literal polynomial
system. Equations \eqref{eq:odd}--\eqref{eq:common-norm} are polynomial
rows of degree at most two in their actual registers.

\subsection{Pell orientation before arithmetic signs}

For a signed integer center \(A\), define the Pell polynomials by
\[
 X_0(A)=1,\quad Y_0(A)=0,\qquad
 X_{k+1}=AX_k+(A^2-1)Y_k,\quad Y_{k+1}=X_k+AY_k.
\]
For negative indices use \(X_{-k}=X_k\), \(Y_{-k}=-Y_k\).
The retained frame bounds every actual base, requested index, and proposed
power output in absolute value by an integer \(W\), with
\begin{equation}\label{eq:frame}
 W\ge258,\qquad |A|>W^W,\qquad |A-1|>4W^3.
\end{equation}
These are consequences of square-based frame equations. Their proof does
not assume signs for \(F,B,q,x\), or the power outputs.

The deterministic comparison center is retained from the supplied
predecessor proof. To distinguish its auxiliary pair
from the arithmetic quantity \(H\), write it as
\[
 H_{\rm p}^2-(A^2-1)J_{\rm p}^2=1,\qquad
 J_{\rm p}=C^2(2e+1),
\]
where \((U,C)\) is the principal anchor pair with requested index \(h\).
The squared cap says
\[
 C^2=h^2+1+\sum_{\text{other pairs}}Y_i^2+G(a',b',c').
\]
It gives \(C\ne0\), \(|h|<|C|\), and strict domination of all other
principal ordinates. The comparison center is defined by
\[
 f=(1-A)H_{\rm p},\qquad D=A+H_{\rm p}f
    =1-(A-1)(A^2-1)J_{\rm p}^2.
\]
Thus \(|D|>1\), \(D\equiv A\pmod{H_{\rm p}}\), and \(D\equiv1\pmod C\).
The comparison ordinate residue is \(R=h+Ck\).

Here is the part of its proof that matters for compatibility. Classify
the anchor and common Pell ordinates at \(|A|\) by positive indices
\(i,l\). For \(i\ge2\), the squared divisibility condition implies
\(|C|\mid l\) and \(l>i\). The abscissa comparison gives the comparison
index \(j\equiv\pm i\pmod{2l}\); the new ordinate residue gives
\(j\equiv\pm h\pmod{|C|}\). Consequently
\(i\equiv\pm h\pmod{|C|}\). The growth bound
\[
 |C|=Y_i(|A|)\ge i+2|A|-2
\]
and \(|h|<W\) imply \(i+|h|<|C|\), so \(i=|h|<W\).
All other classified indices are smaller by the cap.

The exceptional case \(i=1\) has \(|C|=1\). The cap then forces \(h=0\)
and every other ordinate to vanish. Reducing the retained orientation
pool modulo \(A-1\) gives a nonnegative small error sum. Its anchor
ordinate contributes one, so it is strictly positive, while
\eqref{eq:frame} makes it strictly smaller than \(|A-1|\).
Its asserted divisibility is impossible. This exceptional branch
therefore supplies no zero.

All fourteen principal pairs retain both-coordinate orientation:
the squared abscissa errors and squared ordinate-index errors remain
in the pool. After classification their small residues have errors at
most two and strictly less than \(2W\), respectively. The resulting
nonnegative residue sum is less than
\(4\cdot14(W^2+1)<4W^3<|A-1|\). Divisibility forces every error to vanish.
Thus each principal pair is exactly
\((X_k(A),Y_k(A))\) at its signed requested index.

This use of all fourteen pairs is deliberate. A separate predecessor
reduction keeps fewer abscissa tests, but its hypotheses cannot simply
be assumed in the signed bootstrap below. The present construction
does not make that deletion. The complete retained center and frame
proofs are included in the proof closure.

\subsection{A power lemma valid at signed bases}

\begin{lemma}\label{lem:signed-power}
Assume \eqref{eq:frame}. Let \(|n|,|k|,|m|<W\), with \(|n|\ge2\), and set
\[
 T=2An-n^2-1.
\]
If the pair is canonically oriented at \(k\) and its power row asserts
\[
 X_k(A)+(n-A)Y_k(A)\equiv m\pmod T,
\]
then \(k\ge0\) and \(m=n^k\).
For \(n=\epsilon\in\{1,-1\}\), the same row instead gives
\(m=\epsilon^k\) for every integer \(k\), including negative \(k\).
\end{lemma}

\begin{proof}
For \(|n|\ge2\),
\[
 |T|>4W^W-(W-1)^2-1>W^W+W.
\]
Modulo \(T\), one has
\((n-A)^2=A^2-1\) and \(n(2A-n)=1\).
The Pell identities therefore give \(n^k\) at nonnegative indices
and its multiplicative inverse at negative indices. If \(k\ge0\),
the congruence has
\[
 |m-n^k|<W+W^W<|T|,
\]
so it is equality. If \(k=-j<0\), multiplication by \(n^j\) gives
\(mn^j\equiv1\pmod T\). But
\[
 |mn^j-1|\le(W-1)^W+1<|T|,
\]
forcing the impossible integer equality \(mn^j=1\).

For a unit base, \(T=2A\epsilon-2\) and \(2A\equiv2\epsilon\pmod T\).
The sequence
\(f_k=X_k(A)+(\epsilon-A)Y_k(A)\) has initial values
\(f_0=1,f_1=\epsilon\) and satisfies
\(f_{k+1}=2Af_k-f_{k-1}\). This recurrence runs backwards as well
as forwards. Hence \(f_k\equiv\epsilon^k\pmod T\) for every signed
integer \(k\). Finally \(|T|>2W^W-2>W\), whereas
\(|m-\epsilon^k|\le W\), giving the stated equality.
\end{proof}

The base \(-1\) case is essential: it cannot be replaced by a lemma
that presupposes a nonnegative exponent. The literal request for
\(M=q^x\) has the row
\[
 -M+X_x(A)-AY_x(A)+qY_x(A)-T_q\eta=0.
\]
Both \(M\) and \(q\), and the actual base \(BH\) of the self-radix
request, are included in the frame. Thus the lemma applies to the
unshifted arithmetic values claimed here.

